\section{Backgrounds}\label{sec:bakcgournd}
We will begin by covering some basics on quantum error correcting codes (QECC) and their logical operators, as well as describe Kasai's ultra-high rate codes. A glossary of notations is given in Section~\ref{ssec:notations}. 

\subsection{Quantum Error Correcting Codes}\label{ssec:qec-basics}

We say a quantum error correcting code (QECC) is a $[[n,k,d]]$ code if it encodes $k$ logical qubits in the entangled states of $n$ physical qubits, and no Pauli operator of weight less than $d$ maps one logical state to a different logical state. A QECC protects logical information against noise by repeatedly measuring a set of \textit{stabilizers}, which projects the joint state of the physical qubits onto a logical subspace following Pauli corrections determined by the stabilizer measurements. The \emph{encoding rate}, $\frac{k}{n}$, describes the space efficiency in terms of the number of physical qubits needed to encode a logical qubit, and $d$ describes the error correction capacity.


A Pauli operator can be represented as a binary vector $v\in \mathbb F_2^n$, where the $X$-type Pauli $X(v)$ and $Z$-type Pauli $Z(v)$ correspond to
\[
X(v) = \bigotimes_{i\in[n]} X_i^{v_i}, \qquad
Z(v) = \bigotimes_{i\in[n]} Z_i^{v_i},
\]
and the Hamming weight of $v$, $\wt{v} = |\{i: v_i = 1\}|$, corresponds to the weight of the  Pauli operator. Two Pauli operators of different type, i.e.\ $X(v), Z(u)$, commute if and only if their overlap $\rs{v,u} := \sum_i v_iu_i \bmod 2$ vanishes:
\begin{equation}\label{eq:binary-commutation}
    X(v)Z(u) = (-1)^{\rs{v,u}}\, Z(u) X(v).
\end{equation}
 
 
A Calderbank-Shor-Steane (CSS) code is a QECC whose stabilizers are each purely
of $X$- or $Z$-type; by \eqref{eq:binary-commutation}, commutativity reduces to orthogonality of the pair of $X$- and $Z$-stabilizer matrices:

\begin{definition}[CSS Code]\label{def:css}
A CSS code $\mathcal Q = \CSS{H_X, H_Z}$ on $n$ qubits is specified by a pair of binary \emph{stabilizer matrices} $H_X \in \mathbb F_2^{m_X\times n}$ and $H_Z\in \mathbb F_2^{m_Z\times n}$ satisfying the \emph{orthogonality condition}
\begin{equation}\label{eq:css-orthogonality}
    H_X H_Z^T = 0 .
\end{equation}
Each row $x$ of $H_X$ (resp.\ $z$ of $H_Z$) is a Pauli stabilizer $X(x)$ (resp.\ $Z(z)$), and they generate the stabilizer group
\[
\mathcal S = \rs{\, X(x),\, Z(z) \;|\; x = (H_X)_i,\ z = (H_Z)_j,\ 0<i\leq m_X,\ 0<j\leq m_Z\,}.
\]
The logical space of a CSS code is the joint $+1$-eigenspace
$
\mathcal Q = \{\ket \psi : S\ket\psi = \ket \psi \ \forall S \in \mathcal S\},
$
which encodes
\begin{equation}\label{eq:css-k}
    k = n - \rank(H_X) - \rank(H_Z)
\end{equation}
logical qubits.
\end{definition}

The CSS orthogonality condition~\eqref{eq:css-orthogonality} ensures that $\mathcal S$ defines a set of mutually commuting observables, which is necessary for the logical space to be well defined and non-trivial.


Properties of a CSS code can be analyzed using the \emph{Tanner graph} of its check matrices:
\begin{definition}[Tanner Graph Properties]\label{def:qldpc}
Let the \emph{Tanner graph} $\mathcal{T}(H) = G(V_H, E_H)$ of a check matrix $H\in \mathbb F_2^{m\times n}$ be the bipartite graph on vertices $V_H = [n]\sqcup[m]$ with data qubit nodes $[n]$ and check nodes $[m]$, and with edges $E_H = \{(i,j)|H_{i,j} = 1\}$; and let $\mathcal{T}(\mathcal{Q}) =G(V_\mathcal{Q}, E_\mathcal{Q}) =  \mathcal{T}(H_X)\sqcup\mathcal{T}(H_Z)$ denote the Tanner graphs of a CSS code $\mathcal{Q}(H_X,H_Z)$ as in Def.~\ref{def:css}. Then, the CSS code $\mathcal{Q}$ is \emph{qLDPC} if $\deg(v)\leq w$ for all nodes $v\in V_\mathcal{Q}$; or, equivalently, if the row and column weights of both $H_X$ and $H_Z$ are bounded above by a constant $w$. Furthermore, we define the girth of $\mathcal{Q}$ as $\gth{\mathcal Q} = \gth{\mathcal{T}_\mathcal{Q}}$.
\end{definition}
\begin{lemma}[Four-cycle Counts]\label{lemma:girth}
Tanner graphs are bipartite, so girths are even and at least $4$. The number of distinct four-cycles $t_4$ is given by:
\begin{equation}\label{eq:four-cycle-count}
    t_4 = \sum_{0\leq i< j< m_X} \binom{(H_XH_X^T)_{i,j}}{2},
\end{equation}
where the summation is taken over $\mathbb Z$. In particular $\gth{H_X}\geq 6$ if and only if $(H_XH_X^T)_{i,j}\leq 1$ for all $i\neq j$. The same statements hold for $H_Z$.
\end{lemma}
\begin{proof}
Let $c_0,\dots,c_{m_X-1}$ be the check nodes of $\mathcal T(H_X)$ and write $N(c_i) = \supp{(H_X)_i}$. A bipartite graph has no odd cycle and a simple graph has no $2$-cycle, so every cycle length is even and at least $4$. Over $\mathbb Z$,
\begin{equation}\label{eq:gram-overlap}
    (H_XH_X^T)_{i,j} = \sum_{a\in[n]}(H_X)_{i,a}(H_X)_{j,a} = |N(c_i)\cap N(c_j)|,
\end{equation}
since each summand is $1$ exactly when $a$ lies in both supports; over $\mathbb F_2$ only the parity of this overlap would survive.

Map a four-cycle to the pair consisting of its two check indices and its two qubit indices. This is well defined: consecutive vertices of a cycle lie in opposite parts of the bipartition, so a four-cycle alternates and carries exactly two checks and two qubits, say $c_i - q_a - c_j - q_b - c_i$; its vertices are distinct, so $i\neq j$ and $a\neq b$, and all four edges are present, so $a,b\in N(c_i)\cap N(c_j)$. It is injective: the image determines the vertex set $\{c_i,c_j,q_a,q_b\}$, and a four-cycle on that set is unique, because bipartiteness leaves only the four edges $c_iq_a$, $c_iq_b$, $c_jq_a$, $c_jq_b$, and a cycle through all four vertices must give each of them degree $2$ and hence use all four. It is surjective onto the pairs $(\{i,j\},\{a,b\})$ with $i\neq j$, $a\neq b$ and $\{a,b\}\subseteq N(c_i)\cap N(c_j)$, since those edges then form a four-cycle. Grouping by the check pair, each $\{i,j\}$ with $i<j$ admits exactly $\binom{|N(c_i)\cap N(c_j)|}{2}$ choices of $\{a,b\}$, and substituting \eqref{eq:gram-overlap} gives \eqref{eq:four-cycle-count}. Finally the summands are non-negative and $\binom N2 = 0$ iff $N\leq 1$, so $t_4 = 0$ iff every off-diagonal Gram entry is at most $1$.
\end{proof}

\subsubsection{Logical Basis}

Logical qubits of a QECC are usually encoded in entangled $n$-qubit stabilizer states that lack a compact description; therefore, we often use the corresponding logical Pauli operators instead to describe the logical space:

\begin{definition}[Logical Operators and Distance]\label{def:logicals}
Let $\mathcal Q = \CSS{H_X, H_Z}$. The non-trivial logical Pauli operators of $\mathcal{Q}$ are given by 
\begin{align*}
    \mathcal{L}_X = \ns{H_Z}\,/\,\rs{H_X}\;\cong\; \mathbb F_2^{k},\qquad
    \mathcal{L}_Z = \ns{H_X}\,/\,\rs{H_Z}\;\cong\; \mathbb F_2^{k}. 
\end{align*}
The distance of $\mathcal{Q}$ is $d = \min\{d_X, d_Z\}$, where
\begin{equation}\label{eq:css-distance}
    d_X = \min\{\wt{v}: v \in \ns{H_Z}\setminus \rs{H_X}\},\qquad
    d_Z = \min\{\wt{u}: u \in \ns{H_X}\setminus \rs{H_Z}\}.
\end{equation}
\end{definition}
 
It is best to use 
 
\begin{definition}[Orthonormal Logical Basis]\label{def:logical-basis}
A pair of matrices $L_X, L_Z \in \mathbb F_2^{k\times n}$ is a \emph{canonical logical basis} for
$\mathcal Q = \CSS{H_X,H_Z}$ if
\[
L_X H_Z^T = 0,\qquad L_Z H_X^T = 0, \qquad L_XL_Z^T = I_k,
\]
and the rows of $L_X$ (resp.\ $L_Z$) represent a basis of $\ns{H_Z}/\rs{H_X}$ (resp. $\ns{H_X}/\rs{H_Z}$). Such a pair always exists by Gaussian Elimination.
\end{definition}

In classical codes, we can often increase the rate of a code by removing rows of the check matrix to release more logical degrees of freedom, also referred to as a code-modification \emph{augmentation}. However, for quantum codes, the immediate consequence of CSS orthogonality is that any removed stabilizer generator becomes a logical operator itself, unless it was redundant to begin with; hence, if the code we start from is LDPC, the distance after augmentation will be bounded by the same constant $w$:


\begin{fact}[Augmentation Barrier]\label{fact:augmentation}
Let $\mathcal Q = \CSS{H_X,H_Z}$ be a qLDPC CSS code with Tanner-graph degrees $\deg (v)\leq w$ for all $v\in V_{\mathcal Q}$, as in Definition~\ref{def:qldpc}, and let
$\mathcal{Q}^S = \CSS{H^S_X, H^S_Z}$ be an \emph{augmentation} of $\mathcal Q$, obtained by deleting a set $S$ of rows from $H_X$ and from $H_Z$. Suppose some deleted row is not in the row space of the retained rows of its own type; say a deleted row $x$ of $H_X$ has $x\notin \rs{H^S_X}$. Then
\begin{equation}\label{eq:augmentation-barrier}
    d_{\mathcal{Q}^S}\leq d_X(\mathcal Q^S)\leq \wt{x}\leq w.
\end{equation}
The same conclusion holds with the roles of $X$ and $Z$ exchanged.
\end{fact}
\begin{proof}
Since $\mathcal Q$ is a CSS code, $H_XH_Z^T = 0$, and $H^S_Z$ is a submatrix of $H_Z$, so $x (H^S_Z)^T = 0$, i.e.\ $x\in \ns{H^S_Z}$: the operator $X(x)$ commutes with every retained $Z$-stabilizer. By hypothesis $x\notin\rs{H^S_X}$, so $x$ represents a non-trivial class of $\ns{H^S_Z}/\rs{H^S_X} = \mathcal L_X(\mathcal Q^S)$, i.e.\ $X(x)$ is a non-trivial logical $X$ operator of $\mathcal Q^S$. Hence $d_X(\mathcal Q^S)\leq\wt x$, and $\wt x$ is the degree of the check node of $\mathcal T(H_X)$ carrying the row $x$, which is at most $w$.
\end{proof}

Note a linearly independence of the stabilizer is necessary--if the deleted row already lies in the row space of the retained rows of its own type, then $\mathcal Q^S$ and $\mathcal Q$ have the same stabilizer group and hence the same distance. We will assume this to be the case in this work unless otherwise specified.

While not inherently tied to the error correction capacity of a QECC, small cycles in the Tanner graph can lead to failures for  message passing based decoders. We therefore often require $\gth{\mathcal{Q}}\geq 6$ for practical decoder performances, which by Lemma~\ref{lemma:girth} is the condition that every off-diagonal entry of $H_XH_X^T$ and of $H_ZH_Z^T$ be at most $1$ over $\mathbb Z$.


\subsubsection{Chain Complexes}\label{sssec:chain-complex}

We can view a CSS code as length-$2$
chain complex with $H_Z^T,H_X $ as boundary maps with CSS orthogonality corresponding to the boundary map condition:

\begin{definition}[Chain Complex]\label{def:chain-complex}
A length-$\ell$ chain complex over $\mathbb F_2$ is a collection of vector spaces $\{B_i\}_{i=0}^\ell$ together with boundary maps $\partial_i : B_i \to B_{i-1}$ satisfying $\partial_{i-1}\partial_i = 0$, which we abbreviate as $\partial^2 = 0$. A classical code with check matrix $H$ is the length-$1$ complex $C_1 \xrightarrow{H} C_0$, where the bases of $C_1$ and $C_0$ are labeled by bits and checks. A CSS code $\mathcal Q = \CSS{H_X, H_Z}$ is the length-$2$ complex
\begin{equation}\label{eq:css-complex}
    Z \xrightarrow{\ H_Z^T\ } Q \xrightarrow{\ H_X\ } X,
\end{equation}
where the bases of $Z,Q,X$ are labeled by the $Z$-checks, the data qubits, and the $X$-checks respectively, and $\partial_1\partial_2 = H_XH_Z^T = 0$ is exactly~\eqref{eq:css-orthogonality}.

The co-chain complex is the dual:
\begin{equation}\label{eq:css-cocomplex}
    X \xrightarrow{\ H_X^T\ } Q \xrightarrow{\ H_Z\ } Z,
\end{equation}
with coboundary maps $\delta^{i} := \partial_{i+1}^T$; i.e.\ $\delta^0 = H_X^T$ and
$\delta^1 = H_Z$. 
\end{definition}

Logical operators are then the (co)homology of \eqref{eq:css-complex}:
\begin{equation}
    \mathcal L_Z = \ker \partial_1/\operatorname{im}\partial_2 = \ns{H_X}/\rs{H_Z},
    \qquad
    \mathcal L_X = \ker \partial_2^{T}/\operatorname{im}\partial_1^{T} = \ns{H_Z}/\rs{H_X},
\end{equation}

\subsection{Fault Tolerant Logical Operations}\label{ssec:ft-logic}
In the most general definition, a logical operation is the induced transformation of the logical space from a set of physical operations on the data qubits that preserve the logical space; that is, they map an eigenstate to another eigenstate. The logical operation is trivial if it fixes the eigenstates pointwise, and a logical operation is fault-tolerant if failures of these physical operations can be detected and corrected for up to the same distance $d$. An example of a logical operation that is both fault tolerant and trivial is the application of a Pauli stabilizer. We are interested in logical operations that are both non-trivial and fault-tolerant as means to manipulate the encoded logical information. 

There are two main approaches to enable fault tolerant Clifford operations in qLDPC codes: transversal gates, and code deformation. Transversal gates can typically be executed with $O(1)$ time overhead (at least amortized), though their availability depends on the symmetry of the underlying code. While we can always achieve trivial symmetries such as the transversal CNOTs between two identical codes, it is often challenging to attain a useful set of logical gates via happenstance symmetries alone \cite{yang2026spacetimeefficienthardwarecompatiblecomplexquantum}. On the other hand, while code deformation allows us to obtain arbitrary Pauli product measurements by a targeted addition of fault tolerant checks, it typically incurs space-time overheads on the order of $d$ to ensure reliable measurements of the added stabilizers if not carefully optimized. In this work, we mainly investigate gates of the first flavor to prioritize minimizing time overheads, but the symmetries we describe are not specific to that choice: they pay off for surgery- and extractor-based architectures too. 


\subsubsection{Virtual Logical Clifford from Code Automorphisms}
A permutation transversal gate is induced by a permutation of the data qubits that preserves the logical space, also called a code automorphism. They are particularly relevant for RNAA devices, where reconfiguration can be performed with AODs in parallel and with high fidelity. Formally:

\begin{definition}[Code Automorphism]\label{def:code-automorphism}
Let $\mathcal{Q}$ be an $[[n,k,d]]$ CSS code specified by the stabilizers $H_X, H_Z$ and canonical logical basis $L_X, L_Z$. An $n$-dimensional permutation $\sigma$ is said to be a code automorphism of $\mathcal{Q}$ if its right action on the qubit labels preserves the stabilizer group of $\mathcal Q$ sector-wise; that is,
\[
\rs{H_X\sigma} = \rs{H_X},\qquad \rs{H_Z\sigma} = \rs{H_Z},
\]
which we abbreviate as $\mathcal Q\cdot\sigma = \mathcal Q$. Since a permutation preserves overlaps, such a $\sigma$ also maps $\ns{H_Z}$ onto itself, and therefore descends to a linear automorphism of $\mathcal L_X = \ns{H_Z}/\rs{H_X}$ and of $\mathcal L_Z = \ns{H_X}/\rs{H_Z}$. We record the induced logical action as the pair $(g_X, g_Z)$ of $k\times k$ invertible matrices determined by
\[
g_X L_X \equiv L_X \sigma \pmod{\rs{H_X}}, \qquad
g_Z L_Z \equiv L_Z \sigma \pmod{\rs{H_Z}};
\]
the congruences are equalities whenever $\sigma$ happens to fix the chosen representatives, and we then write $g_XL_X = L_X\sigma$.
\end{definition}

We have the following facts:
\begin{fact}[Automorphism Facts]\label{fact:aut-facts}
Let $\mathcal{Q}$ be an $[[n,k,d]]$ CSS code specified by the stabilizers $H_X, H_Z$ and canonical logical basis $L_X, L_Z$. Let $\mathrm{Aut}(\mathcal{Q}) = \{\sigma | \mathcal{Q}\cdot \sigma = \mathcal{Q}\}$ be the set of code automorphisms and let $\mathrm{G}(\mathcal{Q})$ be the set of logical actions $(g_X,g_Z)$ induced by the elements of $\mathrm{Aut}(\mathcal Q)$, as in Definition~\ref{def:code-automorphism}. Then:

\begin{enumerate}
    \item \textbf{Group Homomorphism \cite{berthusen2025automorphismgadgetshomologicalproduct}:} The set of all code automorphism $\mathrm{Aut}(\mathcal{Q})$ and induced logical actions $\mathrm{G}(\mathcal{Q})$ both form a group, and
    \[
    \phi : \mathrm{Aut}(\mathcal{Q}) \to \mathrm{G}(\mathcal{Q}), \qquad \sigma \mapsto (g_X, g_Z)
    \]
    is a group homomorphism.
    \item \textbf{Virtuality:} Suppose $\sigma\in\mathrm{Aut}(\mathcal Q)$ preserves the two stabilizer \emph{matrices} up to a permutation of their rows; that is, suppose there exist row permutations $\rho_X\in S_{m_X}$, $\rho_Z\in S_{m_Z}$ with
    \[
    \rho_X H_X = H_X \sigma, \qquad \rho_Z H_Z = H_Z \sigma.
    \]
    Then the induced logical operation $\phi(\sigma) = (g_X,g_Z)$ is realized by a \emph{virtual relabeling} of the physical qubits: after renaming qubit $i$ as $\sigma(i)$ and renaming the $X$- and $Z$-checks by $\rho_X$ and $\rho_Z$, one recovers the original code together with its original syndrome-extraction schedule, generator by generator. No physical gate is therefore required to implement $\phi(\sigma)$ beyond the rearrangement itself: the gate is a relabeling that can be absorbed into the classical bookkeeping of the decoder, or, on hardware where a fixed physical layout must be restored, into a single atom-rearrangement step.
\end{enumerate}
\end{fact}

Furthermore, we assume that qubit permutation do not spread errors.


\subsubsection{Fold-transversal Logical Clifford from ZX-Duality}
ZX duality also results from a permutation symmetry of the check matrices, but rather than mapping $X$-type stabilizers and $Z$-type stabilizers back to themselves, it maps $X$-type checks to $Z$ and vice versa. Formally,
\begin{definition}[ZX Duality]\label{def:zx-duality}
Let $\mathcal{Q}$ be an $[[n,k,d]]$ CSS code specified by the stabilizers $H_X, H_Z$ and canonical logical basis $L_X, L_Z$, with $m_X = m_Z$. An $n$-dimensional permutation $\tau$ is said to be a ZX duality of $\mathcal{Q}$ if $\tau$ exchanges $H_X$ and $H_Z$ by right action; that is,

\[
H_X\cdot \tau = H_Z,  \quad H_Z\cdot\tau  = H_X.
\]
 $\tau$ is an involution:
\[
H_X\tau^2 = H_Z\tau = H_X,\qquad H_Z\tau^2 = H_X\tau = H_Z.
\]
\end{definition}

Codes with a ZX duality can implement a fold-transversal $H$ gate and phase-type gate, as given in \cite{Breuckmann_2024}:

\begin{fact}[Fold-transversal Gates \cite{Breuckmann_2024}]\label{fact:fold-transversal}
Let $\mathcal Q$ be a CSS code with a ZX duality $\tau$. The following physical operations induce fold-transversal logical operations on $\mathcal Q$:
\begin{align}
    H_\tau &= \bigotimes_{i \in [n],i\leq \tau(i)} \mathrm{SWAP}_{i, \tau(i)} \times \bigotimes_{i \in [n]} H_i,\\
    S_\tau &= \bigotimes_{i \in [n],i\leq \tau(i)} \mathrm{CZ}_{i, \tau(i)} \times \bigotimes_{i \in [n], i = \tau(i)} S_i^\dagger.
\end{align}

\end{fact}


\subsubsection{Homomorphic CNOTs/Measurements from Chain Homomorphisms}
 
The two symmetries above relate a code to itself. Symmetries relating \emph{two} codes also give
rise to transversal gates, and the chain complex formalism of Section~\ref{sssec:chain-complex} is
the natural language for them: a structure-preserving map between two codes is a map between their
complexes that commutes with the boundary maps.
 
\begin{definition}[Chain Homomorphism]\label{def:chain-map}
Let $\mathcal Q, \mathcal Q'$ be CSS codes with boundary maps $\partial_i, \partial'_i$ as in Definition~\ref{def:chain-complex}. A chain homomorphism
$\mathbf \gamma = \{\gamma_z, \gamma_q,\gamma_x\}$ is a triple of linear maps such the following diagram commute:
\begin{equation}\label{eq:chain-map}
    \begin{tikzcd}[ampersand replacement=\&]
        {Z} \& {Q} \& {X} \\
        {Z'} \& {Q'} \& {X'}
        \arrow["{H_Z^T}", from=1-1, to=1-2]
        \arrow["{H_X}", from=1-2, to=1-3]
        \arrow["{H_Z'^T}", from=2-1, to=2-2]
        \arrow["{H_X'}", from=2-2, to=2-3]
        \arrow["{\gamma_z}", from=1-1, to=2-1]
        \arrow["{\gamma_q}", from=1-2, to=2-2]
        \arrow["{\gamma_x}", from=1-3, to=2-3]
    \end{tikzcd}
\end{equation}
That is, $\gamma_z: Z\to Z'$, $\gamma_q: Q\to Q'$ and $\gamma_x: X\to X'$ are named after the space they map \emph{out of}, and commutativity of the two squares reads
\begin{equation}\label{eq:chain-map-squares}
    \gamma_q H_Z^T = H_Z'^T\gamma_z,\qquad
    \gamma_x H_X = H_X'\gamma_q .
\end{equation}
\end{definition}
 
A chain homomorphism preserves the (co)homology of the complex. That is, $\mathbf\gamma$ induces a map $\mathcal L'_Z \to \mathcal L_Z$ on the $Z$-sector and $\mathcal L_X \to \mathcal L'_X$ in the opposite direction. Note that the previous two subsubsections are the special case $\mathcal Q = \mathcal Q'$ with $\gamma_q$ a permutation: a code automorphism is a chain endomorphism of the code, the right-hand square of \eqref{eq:chain-map-squares} becoming $\gamma_x H_X = H_X\gamma_q$, which is the virtuality condition $\rho_XH_X = H_X\sigma$ of Fact~\ref{fact:aut-facts} under $\gamma_x = \rho_X^{-1}$, $\gamma_q = \sigma^{-1}$ (permutations act on the right of check matrices but on the left of chain groups, so the two conventions differ by an inverse, which for a permutation matrix is a transpose). A ZX duality is instead an isomorphism onto the cochain complex \eqref{eq:css-cocomplex}.
 
These homomorphisms identify correspondences between logical qubits of different code blocks, which results in transversal CNOT operators:
 
\begin{definition}[Homomorphic CNOT~\cite{huang2022homomorphiclogicalmeasurements}]\label{def:homomorphic-cnot}
Let $\mathbf \gamma: \mathcal Q \to \mathcal Q'$ be a chain homomorphism as in Definition~\ref{def:chain-map}. The physical $\mathcal Q$-controlled CNOTs specified by $\gamma_q$, i.e.\ a CNOT controlled on qubit $i$ of $\mathcal Q$ and targeting qubit $j$ of $\mathcal Q'$ applied whenever $\gamma_q[i,j] = 1$, implement a set of $\mathcal Q$-controlled logical CNOTs between the two blocks. Here and below we write $\gamma_q$ as the $n\times n'$ incidence matrix of this CNOT pattern, with rows indexed by the qubits of $\mathcal Q$ and columns by those of $\mathcal Q'$, so that it acts on the right of Pauli row vectors; it is the transpose of the vertical arrow of~\eqref{eq:chain-map}. In particular, let $L_X, L_Z$ and $L'_X, L'_Z$ be canonical logical bases for
$\mathcal Q, \mathcal Q'$. The logical circuit is described by
\[
    \prod_{(i,j)\,:\,M_{i,j} = 1} \mathrm{CNOT}(\bar Q_i, \bar Q'_j),
    \qquad\text{where}\qquad
    M = L_X\,\gamma_q\, L'^T_Z \in \mathbb F_2^{k\times k'}.
\]
\end{definition}

The pattern matrix $M$ follows from propagation, since a logical $\bar X_i$ with representative $(L_X)_i$ picks up $(L_X)_i\gamma_q$ on $\mathcal Q'$, whose class is named by pairing against $L'_Z$.
 
We can use these transversal CNOTs to implement logical Pauli product measurements (PPMs). To measure $Z$-type products on a data code $\mathcal Q$, we prepare the logical qubits of an ancilla code $\mathcal Q'$ in the $Z$ basis, apply the $\mathcal Q$-controlled homomorphic CNOTs, and measure $\mathcal Q'$ transversally in the $Z$ basis; the $X$-type version reverses the control and both bases. What gets measured is the set of products named by the \emph{columns} of $M$ --- measuring ancilla logical $j$ returns $\prod_i\bar Z_i^{\,M_{ij}}$ --- so the gadget performs $k'$ measurements in parallel, with single logical qubit measurements being the special case where each column of $M$ has weight one. The scheme is selective, in the sense that the choice of ancilla determines which products are measured.
 
Furthermore, if either block has non-trivial code automorphisms or ZX dualities, the same physical CNOT pattern gives access to a larger family of measurement patterns by composing $\mathbf\gamma$ with the permutations $\sigma$ or $\tau$. Designing the lift so that these symmetries are available by construction is therefore doing double duty, and we
return to it in Section~\ref{sec:logic}.

\subsection{Kasai Codes}\label{ssec:kasai}
Next, we describe the Kasai code construction from
\cite{kasai2026breakingorthogonalitybarrierquantum}. In broad strokes, it generalizes the quasi-cyclic Hagiwara--Imai codes of \cite{Hagiwara_2007}: we start with a
block-circulant proto-matrix $B$ and lift each entry to a $P \times P$ permutation matrix. The key difference is the choice of lift. Whereas the quasi-cyclic construction lifts with commutative cyclic shifts, Kasai lifts with affine permutation matrices (APMs), which in general do \emph{not} commute.

\begin{definition}[Affine Permutation Matrices]
An affine permutation matrix $\mathrm{APM}(a, b, P)$ with $\gcd(a, P) = 1$ is a binary permutation matrix for the map
\[
A: x \mapsto ax + b \mod P.
\]
We use $\mathbb A_P = \{\apmtx{a,b,P}:\gcd(a,P) = 1, 0\leq a,b<P\}$ to denote the group of affine permutations on $[P]$.
\end{definition}

When multiplying two block-circulant matrices, it is helpful to define the following index set
\begin{equation}\label{eq:activeset}
    \Gamma_J= \bigcup_{|r|< J} \left\{(i,j)\, \middle|\, i+j \equiv_{\frac L2} r,\ 0\leq i,j< \frac L2\right\},
\end{equation}
where $r$ ranges over the $2J-1$ residues $-(J-1),\dots, J-1$ modulo $\frac L2$. These are precisely the pairs of generator indices that the first $J$ block rows of the two parent matrices pair against one another: as computed in~\eqref{eq:psi-r} below, the $(i,j)$ block of $\hat H_X\hat H_Z^T$ depends only on the offset $j-i$ and collects the commutators of all generator pairs whose indices sum to $j - i$, and $j-i$ runs over $\{-(J-1),\dots,J-1\}$ as $i,j$ run over $[J]$.

By picking APMs with commutativity relations specified by some active set $\Gamma:\quad \Gamma_{J}\subseteq \Gamma \subseteq [\frac L2]\times [\frac L2]$, we can selectively enforce the CSS orthogonality condition only on a subset of $J$ (active) rows of $B$ while deliberately breaking it on the remaining (latent) rows. This allows the logical degrees of freedom spanned by the latent rows to achieve distances higher than the stabilizer weight. The resulting codes achieve rate $\geq 1-\frac{2J}{L}$, hence $\geq 1/2$ whenever $J\leq \frac L4$, and the check matrices $\hat H_X, \hat H_Z$ obtained from the lift $\{F_i\}, \{G_i\}$ are made explicit below.

\begin{definition}[Kasai Code]\label{def:kasai}
Given parameters $P, L, J \leq \frac{L}{2} \in \mathbb{N}_+$ and an active set $\Gamma:\, \Gamma_{J}\subseteq \Gamma \subseteq [\frac L2]\times [\frac L2]$, the Kasai code $\mathrm K_{L,J}(\mathcal F,\mathcal G)$ is a CSS code specified by generators $\mathcal F = \{F_i\}_{i\in[\frac L2]}, \mathcal{G} = \{G_i\}_{i\in[\frac L2]}\subset \mathbb A_P$ such that $[F_i, G_j] = 0 \iff (i,j)\in \Gamma$. The stabilizers $ H_X, {H}_Z$ are given by the first $J$ rows of the block-circulant parent matrices $\hat{H}_X, \hat{H}_Z$:

\begin{align*}
    [\hat H_X]_{i,j} &= F_{j-i},  &[\hat H_X]_{i,j+\frac L2 } &= G_{j-i}; \\
    [\hat H_Z]_{i,j} &= G^T_{i-j}, & [\hat H_Z]_{i,j+\frac L2 } &= F^T_{i-j}.
\end{align*}
\end{definition}

We call $H_X, H_Z$ the active matrices and $\tilde H_X, \tilde H_Z$ built from the remaining $\frac L2-J$ rows the latent matrices.

To show orthogonality, we compute
\begin{equation}\label{eq:psi-r}
    [\hat H_X\hat H_Z^T]_{i,j} = \Psi_{j-i},\qquad
    \Psi_{r} := \sum_{u\in[\frac L2]} [F_u, G_{r-u}],
\end{equation}
which follows by expanding $[\hat H_X\hat H_Z^T]_{i,j} = \sum_l F_{l-i}G_{j-l} + \sum_l G_{l-i}F_{j-l}$ and substituting $u = l-i$ in the first sum and $u = j-l$ in the second. Then $\Psi_{j-i} = 0$ for each $0\leq i,j<J$, since every term $[F_u, G_{(j-i)-u}]$ has $u + ((j-i)-u) \equiv_{\frac L2} j-i$ with $|j-i|<J$, so the corresponding index pair lies in $\Gamma_J\subseteq\Gamma$ and $[F_i,G_j] = 0\iff (i,j)\in\Gamma$ applies. Hence $H_XH_Z^T = 0$ and the active matrices do define a CSS code.

As a concrete example, take $L = 12$, $J = 3$ and $\Gamma  = \left([\frac L2]\times [\frac L2]\right)  \setminus \{(0,3), (1,2)\}$. We have the $\tfrac L2\times\tfrac L2$ block circulants

\begin{align}
F &= \begin{bmatrix}
F_0 & F_1 & F_2 & F_3 & F_4 & F_5\\
F_5 & F_0 & F_1 & F_2 & F_3 & F_4\\
F_4 & F_5 & F_0 & F_1 & F_2 & F_3\\
F_3 & F_4 & F_5 & F_0 & F_1 & F_2\\
F_2 & F_3 & F_4 & F_5 & F_0 & F_1\\
F_1 & F_2 & F_3 & F_4 & F_5 & F_0
\end{bmatrix},\label{eq:F}\\[6pt]
G &= \begin{bmatrix}
G_0 & G_1 & G_2 & G_3 & G_4 & G_5\\
G_5 & G_0 & G_1 & G_2 & G_3 & G_4\\
G_4 & G_5 & G_0 & G_1 & G_2 & G_3\\
G_3 & G_4 & G_5 & G_0 & G_1 & G_2\\
G_2 & G_3 & G_4 & G_5 & G_0 & G_1\\
G_1 & G_2 & G_3 & G_4 & G_5 & G_0
\end{bmatrix},
\end{align}
from which the parent check matrices are $\hat H_X = [F|G]$ and $\hat H_Z = [G^T|F^T]$,
and 

\begin{align}
\hat H_X \hat H_Z^T 
&=\begin{bmatrix}H_X\\ \tilde H_X\end{bmatrix}
\begin{bmatrix}H_Z^T & \tilde H_Z^T\end{bmatrix}
=\begin{bmatrix}
H_X H_Z^T & H_X\tilde H_Z^T\\
\tilde H_X H_Z^T & \tilde H_X\tilde H_Z^T
\end{bmatrix}\\
&=\left[\begin{array}{ccc|ccc}
\Psi_0 & \Psi_1 & \Psi_2 & \Psi_3 & \Psi_4 & \Psi_5\\
\Psi_5 & \Psi_0 & \Psi_1 & \Psi_2 & \Psi_3 & \Psi_4\\
\Psi_4 & \Psi_5 & \Psi_0 & \Psi_1 & \Psi_2 & \Psi_3\\
\hline
\Psi_3 & \Psi_4 & \Psi_5 & \Psi_0 & \Psi_1 & \Psi_2\\
\Psi_2 & \Psi_3 & \Psi_4 & \Psi_5 & \Psi_0 & \Psi_1\\
\Psi_1 & \Psi_2 & \Psi_3 & \Psi_4 & \Psi_5 & \Psi_0
\end{array}\right] \\
&= 
\left[\begin{array}{ccc|ccc}
&&&\Psi_3&&\\
&&&&\Psi_3&\\
&&&&&\Psi_3\\
\hline
\Psi_3&&&&&\\
&\Psi_3&&&&\\
&&\Psi_3&&&
\end{array}\right]
\end{align}
where, since the only excluded index pairs $(0,3), (1,2)$ both sum to $3$, we have $\Psi_r = 0$ for every $r\neq 3$, and the construction is designed so that $\Psi_3 = [F_0,G_3] + [F_1,G_2]\neq 0$.

A more compact way to write down Kasai's codes is as a two block code where each block is an element from the group algebra of $\mathbb Z_{\frac{L}{2}}\times \mathbb A_P$ acting on $[\frac L2]\times [P]$:

\begin{align}
    F = \sum_{i \in [\frac{L}{2}]} z_i\otimes F_i,\\
    G = \sum_{i \in [\frac{L}{2}]} z_i\otimes G_i;
\end{align}

where $z_i  \in \mathbb Z_{\frac L2}$ is represented as the permutation matrix $[z_{i}]_{j ,k} = \delta_{j+i, k}$.

While achieving high rate, distance, and girth is enough to show strong performance under the code-capacity noise model, practical implementation under circuit level noise imposes additional desiderata, such as shorter code length, hardware compatibility, and explicit, low-weight logical basis and code automorphisms. The follow-up co-design of~\cite{zhao2026ultrahighratequantumerrorcorrection} targets hardware compatibility by relaxing Kasai's girth-$8$ requirement to girth $\geq 6$, and by constraining not the generators themselves but the \emph{transition permutations} of the syndrome-extraction schedule: the rearrangement carrying the ancilla ordering of one round to that of the next is $T_{ij} = F_jF_i^{-1}$ (and likewise $F_jG_i^{-1}$ and $G_jG_i^{-1}$ for the mixed transitions), and one requires every such $T_{ij}$ to commute with a fixed reference APM $A$. Because commuting permutations share an orbit decomposition, each $T_{ij}$ then acts as a cyclic shift within each orbit of $A$ together with a permutation between orbits, i.e.\ as a small number of parallel row and column moves. This yields smaller instances whose syndrome extraction is a short sequence of such structured moves, with millisecond-scale QEC cycle time estimates and low logical error rates (LERs) using only a few pairs of AODs. As we show below, generators subject to this reference-APM constraint factor into a \emph{direct product} or \emph{semi-direct product} of finite groups. We propose instead to lift with product groups directly, retaining the design flexibility of the APM construction while being simpler and substantially more transparent algebraically.

\subsection{Finite Groups and Group Products}\label{ssec:groups}

Let's also establish some notations on finite permutation groups and their matrix representations. 

\begin{definition}[Matrix Permutation Representation]\label{def:mtx-perm}
Given a finite permutation group $G$ acting on $[n]$ by $i\mapsto g(i)$, we define the natural matrix representation as
$$
     M_G : G \to \mathbb F_2^{n\times n},\quad g\mapsto P_g \quad \text{with} \qquad [P_g]_{i,j} = \delta_{j = g(i)}.
$$
We write $M_g$ for $M_G(g)$ when the group is clear from context, and we identify a group element with its matrix throughout. With this convention the matrices act on the right of row vectors, $e_i P_g = e_{g(i)}$, so that $M_gM_h = M_{h\circ g}$.
\end{definition}

\begin{definition}[Direct Product]\label{def:direct-product}
Given finite groups $H_k$ acting on $[k]$ and $C_m$ acting on $[m]$, we can construct a new group via a direct product 
\begin{align*}
     H_k \times C_m = \{(h, c): h\in H_k,\ c\in C_m\}.
\end{align*}
$H_k \times C_m$ acts on the product set $[k]\times[m]$ by $(h,c)\cdot(i,j) = (h(i),\, c(j))$, and this product action is realized by the Kronecker product
\begin{align*}
     M_\times: (h,c) \;\mapsto\; M_H(h)\otimes M_C(c) \;\in\; \mathbb F_2^{km\times km}.
\end{align*}

\end{definition}

\begin{definition}[Semi-direct Product]\label{def:semidirect-product}
Given a finite group $H_k$ acting on $[k]$ and $k$ copies of a finite group $C_m$ acting on $[m]$, we can construct a new group via a semi-direct product (with normal subgroup on the right):
$$
H_k \ltimes C_m^{k} = \{(h;\mathbf c): h\in H_k, \mathbf c =(c_1,...,c_k)\in C_m^k \}
$$
acting on $[km] \cong [k]\times [m]$, such that each element $ (h; \mathbf{c})$ defines the map on $(i, j)\in [k]\times [m]$:
$$
  (h;\mathbf{c}):  (i, j ) \mapsto (h(i), c_{i}(j)).
$$

Suppose $H_k$ and $C_m$ are permutation groups with matrix representations $M_H, M_C$; then the semi-direct product is realized by

\begin{equation}\label{eq:spk-matrix}
  M_\ltimes : (h;\mathbf c)\mapsto
 \left(\bigoplus_{i\in[k]} M_C(c_i)\right)\cdot\left( M_H(h)\otimes I_m\right)
  \;\in\;\mathbb F_2^{km\times km}.
\end{equation}
The order of the two factors matters, and is fixed by the displayed action: reading the right action of Definition~\ref{def:mtx-perm} left to right, $e_{(i,j)}$ is first sent to $e_{(i,c_i(j))}$ by the block-diagonal factor and then to $e_{(h(i),c_i(j))}$ by the block-permutation factor. (The opposite order would instead realize $(i,j)\mapsto (h(i), c_{h(i)}(j))$, which generates the same group but re-labels the bottom vector by $h$.)
\end{definition}


That is, $H_k\ltimes  C_m $ describes permutations on $k$ blocks, each a copy of $[m]$: the $j$-th factor of the base group $C_m^{\,k}$ acts inside block $j$, while $H_k$ permutes between the blocks. Furthermore, if we restrict all factors in $C^k_m$ to be the same, $\mathbf c = (c,\dots,c)$, then $\bigoplus_i M_C(c) = I_k\otimes M_C(c)$ and \eqref{eq:spk-matrix} collapses to $M_H(h)\otimes M_C(c)$; that is, the semidirect product restricts to the direct product $H_k\times C_m$. We also remark that ``product'' here refers to a group product that happens during the lift, as opposed to the code construction technique the word more often refers to, which happens at the proto-matrix level, such as hypergraph-product or balanced-product constructions.

In the following, while this framework generalizes easily to arbitrary groups or matrices, we shall focus only on permutation groups here, and refer to group elements by their matrix representations. Furthermore, the base group $C_m$ is often referred to as the \emph{bottom} of the semi-direct product, while $H_k$ is called the \emph{top}.

We can give some simple facts about the commutativity of these product groups:

\begin{lemma}\label{lemma:dpk-commutivity}
Suppose $H_k$ is non-abelian and $C_m$ is abelian, and let $F = (h,a), G = (h',b)\in H_k\times C_m$. Then
\begin{equation}\label{eq:dpk-commutivity}
    [F,G] = [M_H(h), M_H(h')]\otimes M_C(ab),
    \qquad\text{so}\qquad
    [F,G] = 0 \iff [h,h'] = 0.
\end{equation}
\end{lemma}

\begin{proof}
Using $M_\times(h,c) = M_H(h)\otimes M_C(c)$ from Definition~\ref{def:direct-product}, the mixed-product property of the Kronecker product, and $M_C(a)M_C(b) = M_C(ab) = M_C(ba) = M_C(b)M_C(a)$ (recall $C_m$ is abelian), we have over $\mathbb F_2$
\begin{align*}
        [F,G] &= \left[M_H(h)\otimes M_C(a),\ M_H(h')\otimes M_C(b)\right] \\
        &= \big(M_H(h)\otimes M_C(a)\big) \big( M_H(h')\otimes M_C(b)\big) +\big(M_H(h')\otimes M_C(b)\big)\big(M_H(h)\otimes M_C(a)\big)\\
        &= \big( M_H(h)  M_H(h')\big)\otimes \big(M_C(a) M_C(b)\big) + \big(M_H(h')  M_H(h)\big)\otimes \big(M_C(b) M_C(a)\big)\\
        &= \big( M_H(h)  M_H(h') +  M_H(h')  M_H(h) \big)\otimes M_C(ab)
        = [M_H(h),M_H(h')]\otimes M_C(ab),
\end{align*}
which is the first claim. For the second, $M_C(ab)$ is a permutation matrix, hence non-zero and invertible, and $A\otimes B = 0$ with $B\neq 0$ forces $A = 0$. So $[F,G] = 0$ if and only if $[M_H(h),M_H(h')] = 0$, which is $hh' = h'h$.
\end{proof}

\begin{lemma}\label{lemma:spk-commutivity}
Suppose $H_k$ is non-abelian and $C_m$ is abelian, and let $F=(h;\mathbf a),\ G=(h';\mathbf b)\in H_k\ltimes C_m^{k}$. Then
\begin{equation}\label{eq:spk-commutivity}
[F,G]=0 \iff
\begin{cases}
    [h,h']=0, \quad \text{and} \\
    a_{h'(i)}b_i =b_{h(i)}a_i \quad \forall i\in[k].
\end{cases}
\end{equation}

In particular, there are a few simple one-sided conditions that may be useful for code search:
    \begin{enumerate}
    \item (Necessary) $[F,G]= 0\implies [h,h']= 0$,
    \item (Sufficient) $h=h'=e \implies [F,G]=0$.
    \end{enumerate}
\end{lemma}

\begin{proof}
    Let $A = M_C(a_1)\oplus ...\oplus M_C(a_k)$, $B =M_C(b_1)\oplus ...\oplus M_C(b_k)$, and let $H = M_H(h)\otimes I$, $H' = M_H(h')\otimes I$ we have:
    \begin{align*}
        [F,G] 
        &= [HA , H'B]\\
        &= HAH'B + H'BHA\\
        &= H\cdot(H'{H'}^{T}) \cdot AH' B  + H'\cdot(H  H^T)B HA\\
        &= HH' A_{H'}B + H'H B_HA.
    \end{align*}
    where $A_{H'}={H'}^{T} AH'$ correspond to the bottom vector $(a_{h'(1)},...,a_{h'(k)})$ conjugated by the top $h'$. %It's easy to see that the conjugation acts as a coordinate shift $a_i\mapsto a_{h'(i)}$. 
    (and similarly for $B_H = H^TBH$.)
    %we can see that $a_{h'(i)}b_i = b_{h(i)}a_i\implies A_{H'}B = B_HA$, and $[h,h'] = 0\implies HH' =H'H$, which implies $[F,G] = 0$ together. This shows sufficiency. 

    Let us define $\textbf{p} = (p_1,...,p_k) \in C_m^k $ with $p_i = (a_{h'(i)}b_i a^{-1}_ib^{-1}_{h(i)})$ and $q = (hh')^{-1} h'h\in H_k$ with matrix representations $P = M_{C_m^k}(\textbf{p}) = \bigoplus_{i\in [k]} M_{C_m}(a_{h'(i)}b_i a^{-1}_ib^{-1}_{h(i)})$ and $Q = M_H(q)\otimes I_m$. It's easy to check that $A_{H'}B = P\cdot B_HA$ and $HH'\cdot Q = H'H$ by definition. Finally, we have $ [F,G] = HH' A_{H'}B + H'H B_HA = HH'(P+Q) B_HA$. Since $B_HA \neq 0$ and $HH'\neq 0$, and $[F,G] = 0\iff P = Q = I$, which is equivalent to the conditions in ~\ref{eq:spk-commutivity}. The rest of the simple conditions follow from straight forwardly. 
    
\end{proof}

\subsection{Notations}\label{ssec:notations}
Finally, we give a glossary of notations used throughout the supplementary materials:
\begin{enumerate}
    \item $[J] = \{0,1,...,J-1\}$ denotes the first $J$ natural numbers.
    \item $S_n$ denotes the permutation group acting on $[n]$.
    \item $\mathbb Z_n$ denotes the cyclic group acting on $[n]$.
    \item $\mathbb A_n$ denotes the affine group acting on $[n]$.
    \item $[A, B] = AB-BA$ denotes the commutator; over $\mathbb F_2$ this is $AB+BA$. For group elements $[f,g] = 0$ means $[M_f,M_g] = 0$, equivalently $fg = gf$.
    \item $a \equiv_P b$ denotes $a = b\mod P$.
    \item $A\oplus B = \begin{bmatrix}A&\\&B\end{bmatrix}$ denotes the direct sum of two matrices $A,B$.
    \item $\mathbb F^{n\times m}_q$ denotes the space of $n\times m$ dimensional $q$-ary matrices.
    \item $\textbf{1}^{n\times m} \in \mathbb F_2^{n\times m}$ denotes the $n\times m$ matrix with all $1$ entries.
    \item $Z_G(S) = \{A\in G, [A,s] = 0\ \forall s\in S\}$ denotes the centralizer of $S$ in $G$. $Z_G(S)$ fixes $S$ point-wise under conjugation.
    \item $N_G(S) = \{A\in G, AsA^{-1} \in S\,  \forall s\in S\}$ denotes the normalizer of $S$ in $G$. $N_G(S)$ fixes $S$ set-wise.
    \item $\pr$ denotes an \emph{orbit projector}. For a product group acting on $[k]\times[m]\cong[n]$, $n = km$, as in Definitions~\ref{def:direct-product} and~\ref{def:semidirect-product}, we write
    \[
        \pr_C = I_k\otimes \textbf{1}^{1\times m}\in \mathbb F_2^{k\times n},
        \qquad
        \pr^H = \textbf{1}^{1\times k}\otimes I_m\in \mathbb F_2^{m\times n},
    \]
    for the maps that sum the coordinates over each orbit of the bottom factor $C_m$, respectively of the top factor $H_k$. In both cases the decoration records the factor that is \emph{collapsed}, so $\pr_C$ retains the top index and $\pr^H$ the bottom index. Their defining property is that they intertwine the lift with its factors: for the direct product,
    \[
        \pr_C\, M_\times(h,c) = M_H(h)\,\pr_C,
        \qquad
        \pr^H\, M_\times(h,c) = M_C(c)\,\pr^H,
    \]
    which is what makes the quotient codes of Definition~\ref{def:quotient} well defined. More generally, for a subgroup $N\leq G$ all of whose orbits on $[n]$ have the same size, $\pr_N$ denotes the corresponding orbit-sum matrix, with one row per orbit. Note that the $n\times n$ orbit sum $\sum_{i\in[m]}M_C(c)^i = I_k\otimes\textbf{1}^{m\times m}$, taken over a generator $c$ of a cyclic $C_m$ acting with $k$ full orbits, has the same row space $\rs{\pr_C}$ and may be used interchangeably where only the row space matters.
\end{enumerate}